% GENERATED FILE. Edit canonical lessons, then run npm run generate:books.
% canonical-source-hash: fnv1a64:61fb4aee03bf0399
% canonical-lessons: BN-C392-traphik, BN-C392-pothchari, BN-C392-phlaiobhar, BN-C392-sainbord, BN-C392-ghesha

\chapter{Traffic and Crossings}
\label{ch:bn-traffic-and-crossings}
\hlchaptermodality{392}

\begin{chapteropening}
\textbf{By the end of this chapter:} \emph{I can find my way through traffic, past signboards and over a flyover.}

Complete the last lesson of chapter 392: I can find my way through traffic, past signboards and over a flyover.
\end{chapteropening}

\section[\texorpdfstring{\bn{ট্রাফিক} (ṭrāphik)}{ট্রাফিক (ṭrāphik)}]{\bn{ট্রাফিক} (ṭrāphik) — traffic (on the roads)}

\label{lesson:BN-C392-traphik}



\begin{quote}
Before the new one: say the Bengali for to overtake, then the Bengali for to apply (the brakes).
\end{quote}

\subsection*{You'll want to know: \bn{ট্রাফিক}}
\textbf{\bn{ট্রাফিক}} — \emph{ṭrāphik} — \textquotedblleft{}traffic (on the roads)\textquotedblright{}.

\begin{quote}
\bn{ট্রাফিক} \bn{পুলিশ} \bn{হাত} \bn{তুলে} \bn{গাড়ি} \bn{থামালেন।}
\end{quote}

\begin{quote}
\emph{ṭrāphik pulish hāt tule gāṛi thāmālen.}
\end{quote}

\begin{quote}
The traffic policeman raised his hand and stopped the cars.
\end{quote}

\bn{যানজট} is a jam; \bn{ট্রাফিক} is the traffic itself, and \bn{ট্রাফিক} \bn{পুলিশ} the officer at the crossing.

\begin{cousinweb}
\textbf{\bn{ট্রাফিক}} (\emph{ṭrāphik}) is borrowed from English \emph{traffic}; Bengali has no f, so the English f is said with \bn{ফ}.
\end{cousinweb}

\begin{grammarlens}[title={What you've built}]
One more word for getting around.
\end{grammarlens}

\subsection*{Your turn}
\begin{itemize}
\raggedright
  \item \emph{Say it:} \emph{ṭrāphik}
  \item \emph{Say it:} \emph{ṭrāphik}, once more
  \item \emph{Say it:} \emph{ṭrāphik pulish hāt tule gāṛi thāmālen.} — The traffic policeman raised his hand and stopped the cars.
  \item \emph{Recall:} say the Bengali for to overtake, then the Bengali for to apply (the brakes), then say \emph{ṭrāphik} again
\end{itemize}

\begin{tcolorbox}[breakable,colback=teal!4,colframe=teal!35!black,title={Before you move on}]
What does \bn{ট্রাফিক} mean? (\textquotedblleft{}Traffic (on the roads)\textquotedblright{}.) Say it once more.
\end{tcolorbox}
\section[\texorpdfstring{\bn{পথচারী} (pôthchārī)}{পথচারী (pôthchārī)}]{\bn{পথচারী} (pôthchārī) — a pedestrian}

\label{lesson:BN-C392-pothchari}



\begin{quote}
Before the new one: say the Bengali for to apply (the brakes), then the Bengali for traffic.
\end{quote}

\subsection*{You'll want to know: \bn{পথচারী}}
\textbf{\bn{পথচারী}} — \emph{pôthchārī} — \textquotedblleft{}a pedestrian\textquotedblright{}.

\begin{quote}
\bn{পথচারীরা} \bn{মোড়ে} \bn{রাস্তা} \bn{পার} \bn{হল।}
\end{quote}

\begin{quote}
\emph{pôthchārīrā moṛe rāstā pār hôlô.}
\end{quote}

\begin{quote}
The pedestrians crossed the road at the corner.
\end{quote}

\begin{cousinweb}
\textbf{\bn{পথচারী}} (\emph{pôthchārī}) is built from \textbf{\bn{পথ}}, \textquotedblleft{}a way\textquotedblright{}, and \emph{-cārin}, \textquotedblleft{}moving, going about\textquotedblright{} (Monier-Williams for both): one who goes on foot along the way.
\end{cousinweb}

\begin{grammarlens}[title={What you've built}]
One more word for getting around.
\end{grammarlens}

\subsection*{Your turn}
\begin{itemize}
\raggedright
  \item \emph{Say it:} \emph{pôthchārī}
  \item \emph{Say it:} \emph{pôthchārī}, once more
  \item \emph{Say it:} \emph{pôthchārīrā moṛe rāstā pār hôlô.} — The pedestrians crossed the road at the corner.
  \item \emph{Recall:} say the Bengali for to apply (the brakes), then the Bengali for traffic, then say \emph{pôthchārī} again
\end{itemize}

\begin{tcolorbox}[breakable,colback=teal!4,colframe=teal!35!black,title={Before you move on}]
What does \bn{পথচারী} mean? (\textquotedblleft{}A pedestrian\textquotedblright{}.) Say it once more.
\end{tcolorbox}
\section[\texorpdfstring{\bn{ফ্লাইওভার} (phlāiobhār)}{ফ্লাইওভার (phlāiobhār)}]{\bn{ফ্লাইওভার} (phlāiobhār) — a flyover}

\label{lesson:BN-C392-phlaiobhar}



\begin{quote}
Before the new one: say the Bengali for traffic, then the Bengali for a pedestrian.
\end{quote}

\subsection*{You'll want to know: \bn{ফ্লাইওভার}}
\textbf{\bn{ফ্লাইওভার}} — \emph{phlāiobhār} — \textquotedblleft{}a flyover\textquotedblright{}.

\begin{quote}
\bn{নতুন} \bn{ফ্লাইওভারে} \bn{যানজট} \bn{নেই।}
\end{quote}

\begin{quote}
\emph{nôtun phlāiobhāre jānjôṭ nei.}
\end{quote}

\begin{quote}
There is no traffic jam on the new flyover.
\end{quote}

\begin{cousinweb}
\textbf{\bn{ফ্লাইওভার}} (\emph{phlāiobhār}) is borrowed from English \emph{flyover}, with \bn{ফ} for f and \bn{ভ} for v.
\end{cousinweb}

\begin{grammarlens}[title={What you've built}]
One more word for getting around.
\end{grammarlens}

\subsection*{Your turn}
\begin{itemize}
\raggedright
  \item \emph{Say it:} \emph{phlāiobhār}
  \item \emph{Say it:} \emph{phlāiobhār}, once more
  \item \emph{Say it:} \emph{nôtun phlāiobhāre jānjôṭ nei.} — There is no traffic jam on the new flyover.
  \item \emph{Recall:} say the Bengali for traffic, then the Bengali for a pedestrian, then say \emph{phlāiobhār} again
\end{itemize}

\begin{tcolorbox}[breakable,colback=teal!4,colframe=teal!35!black,title={Before you move on}]
What does \bn{ফ্লাইওভার} mean? (\textquotedblleft{}A flyover\textquotedblright{}.) Say it once more.
\end{tcolorbox}
\section[\texorpdfstring{\bn{সাইনবোর্ড} (sāinborḍ)}{সাইনবোর্ড (sāinborḍ)}]{\bn{সাইনবোর্ড} (sāinborḍ) — a signboard}

\label{lesson:BN-C392-sainbord}



\begin{quote}
Before the new one: say the Bengali for a pedestrian, then the Bengali for a flyover.
\end{quote}

\subsection*{You'll want to know: \bn{সাইনবোর্ড}}
\textbf{\bn{সাইনবোর্ড}} — \emph{sāinborḍ} — \textquotedblleft{}a signboard\textquotedblright{}.

\begin{quote}
\bn{সাইনবোর্ডটা} \bn{দেখে} \bn{আমরা} \bn{ডানদিকে} \bn{গেলাম।}
\end{quote}

\begin{quote}
\emph{sāinborḍṭā dekhe āmrā ḍāndike gelām.}
\end{quote}

\begin{quote}
We saw the signboard and turned right.
\end{quote}

\begin{cousinweb}
\textbf{\bn{সাইনবোর্ড}} (\emph{sāinborḍ}) is borrowed from English \emph{signboard}, the painted board over a shop or at a turning.
\end{cousinweb}

\begin{grammarlens}[title={What you've built}]
One more word for getting around.
\end{grammarlens}

\subsection*{Your turn}
\begin{itemize}
\raggedright
  \item \emph{Say it:} \emph{sāinborḍ}
  \item \emph{Say it:} \emph{sāinborḍ}, once more
  \item \emph{Say it:} \emph{sāinborḍṭā dekhe āmrā ḍāndike gelām.} — We saw the signboard and turned right.
  \item \emph{Recall:} say the Bengali for a pedestrian, then the Bengali for a flyover, then say \emph{sāinborḍ} again
\end{itemize}

\begin{tcolorbox}[breakable,colback=teal!4,colframe=teal!35!black,title={Before you move on}]
What does \bn{সাইনবোর্ড} mean? (\textquotedblleft{}A signboard\textquotedblright{}.) Say it once more.
\end{tcolorbox}
\section[\texorpdfstring{\bn{ঘেঁষা} (ghẽshā)}{ঘেঁষা (ghẽshā)}]{\bn{ঘেঁষা} (ghẽshā) — to keep close to; to brush against}

\label{lesson:BN-C392-ghesha}



\begin{quote}
Before the new one: say the Bengali for a flyover, then the Bengali for a signboard.
\end{quote}

\subsection*{You'll want to know: \bn{ঘেঁষা}}
\textbf{\bn{ঘেঁষা}} — \emph{ghẽshā} — \textquotedblleft{}to keep close to; to brush against\textquotedblright{}.

\begin{quote}
\bn{সাইকেলটা} \bn{ফুটপাত} \bn{ঘেঁষে} \bn{চলল।}
\end{quote}

\begin{quote}
\emph{sāikelṭā phuṭpāt ghẽshe chôllô.}
\end{quote}

\begin{quote}
The bicycle kept close to the pavement.
\end{quote}

Past (he/she …): \bn{সেই} \bn{গাড়ি} \bn{দেয়াল} \bn{ঘেঁষল}, \emph{sei gāṛi deyāl ghẽshlô}.

Most often in the -\bn{এ} form before a verb of motion: \bn{দেয়াল} \bn{ঘেঁষে} \bn{দাঁড়ানো}, to stand close against the wall.

\begin{cousinweb}
\textbf{\bn{ঘেঁষা}} (\emph{ghẽshā}) is inherited: CDIAL lists Bengali \emph{ghẽs}, \textquotedblleft{}a touch, a rub\textquotedblright{}, under Sanskrit \emph{gharṣa}, \textquotedblleft{}friction\textquotedblright{} (CDIAL 4448); to brush past is to touch in passing.
\end{cousinweb}

\begin{grammarlens}[title={What you've built}]
One more word for getting around.
\end{grammarlens}

\subsection*{Your turn}
\begin{itemize}
\raggedright
  \item \emph{Say it:} \emph{ghẽshā}
  \item \emph{Say it:} \emph{ghẽshā}, once more
  \item \emph{Say it:} \emph{sāikelṭā phuṭpāt ghẽshe chôllô.} — The bicycle kept close to the pavement.
  \item \emph{Recall:} say the Bengali for a flyover, then the Bengali for a signboard, then say \emph{ghẽshā} again
\end{itemize}

\begin{tcolorbox}[breakable,colback=teal!4,colframe=teal!35!black,title={Before you move on}]
What does \bn{ঘেঁষা} mean? (\textquotedblleft{}To keep close to; to brush against\textquotedblright{}.) Say it once more.
\end{tcolorbox}
